\section{Introduction}
\label{sec:introduction}

From climate change and critical minerals to energy resources and disaster risk, Earth science seeks to derive testable scientific conclusions from observational data. These research frontiers are data intensive, computationally demanding, and interdisciplinary. Geoscientific research is shifting from traditional empirical inference and mechanistic simulation toward predictive systems science driven jointly by big data and artificial intelligence~\cite{zheng2025frontiers}. In this emerging paradigm, research capability depends not only on the scale of data and models, but also on the ability to organize distributed data, tools, and methods into reproducible, traceable analytical workflows.

Traditional geoscientific analysis relies heavily on researchers manually connecting data-processing stages. Data are distributed across databases, publications, and repositories, while differences in reporting standards impede comparison and integration~\cite{zheng2025frontiers}. Researchers typically locate data, align them spatially or temporally, derive variables, perform statistical modeling, and verify the results. Resource and disaster analyses may additionally combine data retrieval, spatial processing, and interdisciplinary models. This process is time consuming, labor intensive, difficult to scale, and difficult to reproduce. Scientific agents built around large language models (LLMs) offer an opportunity to automate these workflows: they translate natural-language tasks into multistep execution plans, coordinate multiple tools, and have shown considerable promise in scientific analysis~\cite{feng2026earthagent,ding2025scitoolagent}. Geoscientific workflows, however, have a defining feature that distinguishes them from ordinary question answering: later operations depend on artifacts produced by earlier steps. Model memory or a single tool call is insufficient to complete the entire processing chain. The value of scientific agents for geoscientific research therefore depends on their ability to reliably orchestrate multistep tool chains and produce conclusions grounded in execution.

Existing geoscientific agents have begun to support multistep tool execution using established frameworks such as ReAct~\cite{yao2023react}. For example, GeoAgent~\cite{chen2024geoagent} and Earth-Agent~\cite{feng2026earthagent} use planning--execution--reflection architectures to process Earth science data through multiple calls. These studies establish the feasibility of multistep tool use, but important limitations remain when such agents are expected to orchestrate complete analytical workflows. In complex workflows, their basic execution units are typically individual calls, individual tools, or entire Skills. Planning directly around tool names couples the plan to changes in tool interfaces. Reusing an entire Skill, in turn, restricts reuse to predefined workflows and prevents task-specific extraction and recombination of local operations from different strategies. Operation-level reuse and preservation of usage constraints remain unresolved in geoscientific settings. An open question is how to select and combine the capabilities required for a task at a finer operational granularity, preserve their prerequisites and usage constraints, and match them to suitable concrete tools.

A second challenge is the need for a benchmark that can assess tool orchestration by geoscientific agents. Evaluation has evolved from testing knowledge to examining execution processes. Early work centered on knowledge questions and multiround scientific conversations. EarthSE, for example, evaluates geoscientific reasoning through question answering and exploratory dialogue~\cite{xu2025earthse}, but does not require actual tool execution and therefore cannot assess the analytical process itself. More recently, GeoBenchX~\cite{krechetova2025geobenchx}, Earth-Bench~\cite{feng2026earthagent}, and GeoAgentBench~\cite{yu2026geoagentbench} explicitly examine multistep tool use and execution; the latter two also evaluate processes and final outcomes. We further examine tool composition and orchestration across geoscientific themes. Within the same task, agreement with the reference answer and completion of the required operations should be assessed separately. Failure analysis additionally requires parameters, artifacts, and invocation records, rather than equating answer agreement or an individual successful call with a correct analytical process.

To address these two challenges, we introduce AgentGS, a tool-orchestration method based on hierarchical capability composition, and GeoToolEval, a benchmark for geoscientific tool-workflow orchestration. AgentGS provides a method for acquiring orchestration capabilities, while GeoToolEval provides an evaluation setting for complete geoscientific tool chains. Our main contributions are as follows:

\begin{itemize}
\item We propose AgentGS, a geoscientific tool-orchestration agent using a Skill--Capability--Tool hierarchy and hypergraph-based missing-capability completion. It supports reuse across strategies and constrained tool selection, connecting capability completion, parameter and artifact binding, and local recovery into an executable workflow.
\item We construct GeoToolEval with 603 evaluation records spanning ten geoscientific research-frontier themes. Single-step and multistep tasks associate final-answer evaluation with tool composition, invocation order, and execution-record analysis within the same task, supporting analysis of orchestration deviations.
\item AgentGS achieves a mean FAS of 0.8731 on all 603 GeoToolEval records, exceeding the evaluated language-model and agent baselines. Multistep strata and configuration comparisons reveal differences in answer scores and tool-workflow consistency; their causes require analysis of actual execution traces.
\end{itemize}
